\appendix
\section{Archivos incluidos}

El proyecto entregado incluye los siguientes directorios:

\begin{itemize}
    \item \texttt{tex/}: secciones del informe en \LaTeX.
    \item \texttt{figuras/}: gráficos generados y capturas utilizadas en el informe.
    \item \texttt{data/}: archivos CSV medidos y simulados.
    \item \texttt{ltspice/}: archivos \texttt{.asc} de LTspice.
    \item \texttt{codigo/}: código de la interfaz gráfica en Python.
\end{itemize}

\section{Archivos de LTspice relevantes}

\begin{table}[H]
    \centering
    \begin{tabular}{ll}
        \toprule
        Archivo & Descripción \\
        \midrule
        \texttt{opcionA.asc} & Alternativa mixta RLC y RC con buffer. \\
        \texttt{opcionB.asc} & Alternativa con dos etapas RLC y buffer. \\
        \texttt{filtroFINAL.asc} & Filtro RC final implementado para PCB. \\
        \texttt{2\_1.asc} & Simulación transitoria. \\
        \texttt{2\_2.asc} & Simulación de respuesta en frecuencia. \\
        \bottomrule
    \end{tabular}
    \caption{Archivos de simulación incluidos en el proyecto.}
\end{table}

\section{Código de la interfaz}

El código completo de la interfaz se encuentra en:

\begin{center}
    \texttt{codigo/gui\_tc1\_user\_friendly\_pdf.py}
\end{center}

La interfaz puede ejecutarse con:

\begin{lstlisting}[style=pythonstyle]
py codigo/gui_tc1_user_friendly_pdf.py
\end{lstlisting}

\section{Datos experimentales y simulados}

Los archivos CSV utilizados para generar los gráficos del informe se encuentran en la carpeta \texttt{data/}. En particular, se incluyeron datos de transitorios, respuesta en frecuencia de circuitos de segundo orden y respuesta en frecuencia del filtro final medido en protoboard y en placa.
